% Paper: Galois orbit structure of rational CM weight-5 newforms
% Target: Mathematics of Computation
% Author: Kevin Remondiere
% Date: 2026-05-09

\documentclass[12pt]{amsart}

\usepackage{amsmath, amssymb, amsthm}
\usepackage{hyperref}
\usepackage{booktabs}
\usepackage{microtype}

% Theorem environments
\newtheorem{theorem}{Theorem}[section]
\newtheorem{proposition}[theorem]{Proposition}
\newtheorem{lemma}[theorem]{Lemma}
\newtheorem{corollary}[theorem]{Corollary}
\theoremstyle{definition}
\newtheorem{definition}[theorem]{Definition}
\newtheorem{example}[theorem]{Example}
\newtheorem{remark}[theorem]{Remark}

% Shorthand
\newcommand{\Q}{\mathbb{Q}}
\newcommand{\Z}{\mathbb{Z}}
\newcommand{\C}{\mathbb{C}}
\newcommand{\N}{\mathbb{N}}
\newcommand{\OK}{\mathcal{O}_K}
\newcommand{\ClK}{\mathrm{Cl}(K)}
\newcommand{\Gal}{\mathrm{Gal}}
\newcommand{\HK}{K^{\mathrm{H}}}
\newcommand{\Nm}{\mathrm{N}}
\newcommand{\psi}{\psi}
\newcommand{\alphatwo}{\alpha_2}

\begin{document}

\title{Galois orbit structure of rational CM weight-5 newforms
  over imaginary quadratic fields of class number $h$}

\author{Kevin Remondiere}

\date{May 2026}

\begin{abstract}
Let $K = \Q(\sqrt{D})$ be an imaginary quadratic field of class number $h(K)$,
and let $f$ be a CM weight-5 newform attached to a Hecke Gr\"ossencharacter
of $K$ with infinity type $(4, 0)$.
We prove that the rational-over-$\Q$ CM newforms at conductor $\mathfrak{N}(\OK)$
form a single orbit under the action of $\ClK$ on Gr\"ossencharacters by
ideal-class twisting, and that this orbit has cardinality exactly $h(K)$.
In particular: the eight Heegner (class-number-1) fields each yield one
rational eigenform (or a twin pair); the two illustrative class-number-2 fields
$D \in \{-88, -91\}$ each yield exactly two rational eigenforms related by a
Galois sign-flip; and the class-number-4 field $D = -84$ yields exactly four
rational eigenforms, verified computationally via PARI/GP.
A key invariant of the orbit is that $a_5 = 16$ for all members whenever
$5$ is inert in $K$ and $D$ is a Heegner-like discriminant, matching the
pattern found at $D \in \{-19, -43, -67, -163\}$.
We further clarify a distinction critical for applications:
a \emph{rational eigenform} ($a_p \in \Q$ for all $p$) is
\emph{not} the same as a \emph{rational $L$-value quotient}
$\alphatwo = L(f, 2)\pi^2/\Omega^4 \in \Q$,
which holds only for $h(K) = 1$ (by a theorem of Goldstein--Schappacher).
The case $h(K) = 3$ yields zero rational eigenforms, explained by the
non-triviality of the class-group action on the infinity type.
\end{abstract}

\maketitle

\tableofcontents

% ============================================================
\section{Introduction}
% ============================================================

\subsection{Classical background}
The arithmetic of CM modular forms is governed by Hecke Gr\"ossencharacters
of imaginary quadratic fields.  Shimura \cite{Shimura1971} established that
every weight-$k$ CM newform $f$ of conductor $N$ arises from a Gr\"ossencharacter
$\psi : \mathbb{A}_K^\times \to \C^\times$ with infinity type $(\ell, 0)$,
$\ell = k-1$, via
\[
  f = \theta(\psi) := \sum_{\mathfrak{a} \subset \OK} \psi(\mathfrak{a}) q^{\Nm\mathfrak{a}}.
\]
The Galois group $\Gal(\HK/K)$ of the Hilbert class field of $K$ acts on
the set of such characters by twisting:
\[
  \psi^\sigma(\mathfrak{a}) := \sigma(\psi(\sigma^{-1}\mathfrak{a})), \quad \sigma \in \Gal(\HK/\Q).
\]
This action permutes the associated newforms, producing what we call the
\emph{$\ClK$-orbit} of $f$.

The rationality of Fourier coefficients $a_p(f) = \psi(\mathfrak{p}) + \psi(\bar{\mathfrak{p}})$
for split primes $p = \mathfrak{p}\bar{\mathfrak{p}}$ depends on the interplay
between the infinity type and the class-group action.
For split $p$ the coefficient $a_p = \psi(\mathfrak{p}) + \overline{\psi(\mathfrak{p})}$
lies in $\Q$ if and only if $\psi(\mathfrak{p}) \in \R$, which for
infinity type $(\ell, 0)$ and $\ell$ even is automatic.
For inert $p$ one has $a_p = 0$ by definition.
The combined effect is that \emph{all} Fourier coefficients of $\theta(\psi)$
lie in $\Q$ precisely when the character values at split ideals are real,
a condition that Goldstein and Schappacher \cite{GoldsteinSchappacher1981}
analysed systematically via the theory of CM types.

\subsection{Motivation and new content}
The question of \emph{counting} the number of rational CM newforms of a given
weight at a given level has not, to our knowledge, been explicitly compiled
in the literature, although the ingredients are classical.
The present paper does exactly this for weight~5 (infinity type $(4,0)$),
verifies the count computationally for several discriminants, and
records two conceptual distinctions that are important for applications:

\begin{enumerate}
\item[(i)] The count of rational eigenforms equals $h(K)$, not a smaller number;
  in particular it grows with the class number.
\item[(ii)] Rationality of $a_p$ does \emph{not} imply rationality of the
  special $L$-value quotient $\alphatwo = L(f, 2)\pi^2/\Omega^4$;
  the latter is rational only for $h(K) = 1$.
\end{enumerate}

Point (ii) recovers Goldstein--Schappacher \cite[Th\'eor\`eme~5.1]{GoldsteinSchappacher1981}
from a different angle, and we record it here for clarity.

\subsection{Organisation}
Section~\ref{sec:setup} fixes notation.
Section~\ref{sec:theorem} states and proves the main theorem.
Section~\ref{sec:verification} gives the verification table with PARI/GP data.
Section~\ref{sec:distinction} discusses the distinction between rational eigenforms
and rational $L$-value quotients.
Section~\ref{sec:open} discusses the anomalous $h = 3$ case and open questions.

% ============================================================
\section{Setup and notation}
\label{sec:setup}
% ============================================================

\subsection{Imaginary quadratic fields and class groups}

Let $D < 0$ be a fundamental discriminant, $K = \Q(\sqrt{D})$,
$\OK$ its ring of integers, $h = h(K) = |\ClK|$ the class number,
and $\HK$ the Hilbert class field of $K$.
Write $\chi_D = \left(\frac{D}{\cdot}\right)$ for the Kronecker symbol.
For a rational prime $p$, we say $p$ is \emph{split}, \emph{inert},
or \emph{ramified} in $K$ according as $\chi_D(p) = +1, -1, 0$ respectively.

\subsection{Hecke Gr\"ossencharacters of infinity type $(4,0)$}

A Hecke Gr\"ossencharacter $\psi$ of $K$ with conductor $\mathfrak{f}$
and infinity type $(4,0)$ is a continuous homomorphism
\[
  \psi : \mathbb{A}_K^\times \longrightarrow \C^\times
\]
such that $\psi(x_\infty) = x_\infty^4$ for $x_\infty \in K \otimes_\Q \R \cong \C^\times$,
and $\psi$ is trivial on $(\hat{\OK})^\times \cap \ker(\mathfrak{f})$.
For $K$ of class number $h$ and trivial conductor $\mathfrak{f} = \OK$
the number of such characters is exactly $h$:
they are parametrised by the $h$ choices of a fractional ideal
representative $[\mathfrak{a}] \in \ClK$ used to ``twist'' the trivial character.

More precisely, fix one solution $\psi_0$ to the infinity-type condition
(which exists and is unique up to finite-order twist when $h = 1$).
For general $h$, the full set of solutions is
\begin{equation}
  \label{eq:orbit}
  \{ \psi_0 \otimes \chi_{\mathfrak{a}} : [\mathfrak{a}] \in \ClK \},
\end{equation}
where $\chi_{\mathfrak{a}}$ is the unramified character of $\mathbb{A}_K^\times$
associated to $[\mathfrak{a}]$ via class-field theory.
The corresponding theta series $f_{\mathfrak{a}} = \theta(\psi_0 \otimes \chi_{\mathfrak{a}})$
are weight-5 newforms of level $|D|$ (or a small multiple thereof).

\subsection{The Galois sign-flip action}

Complex conjugation $\tau \in \Gal(K/\Q)$ acts on $\psi$ by
$\psi^\tau(\mathfrak{a}) = \overline{\psi(\bar{\mathfrak{a}})}$.
For infinity type $(4, 0)$ this gives $\psi^\tau = \overline{\psi}$
(the contragredient character), and one checks that
\[
  a_p(f_{\mathfrak{a}}^\tau) = \overline{a_p(f_{\mathfrak{a}})} = a_p(f_{\mathfrak{a}})
\]
for all $p$ inert in $K$, and
\[
  a_p(f_{\mathfrak{a}}^\tau) = \psi^\tau(\mathfrak{p}) + \psi^\tau(\bar{\mathfrak{p}})
    = \overline{\psi(\bar{\mathfrak{p}})} + \overline{\psi(\mathfrak{p})}
    = \overline{a_p(f_{\mathfrak{a}})}
\]
for split $p$.  Since $a_p \in \Q$ by hypothesis (infinity type even),
$\overline{a_p} = a_p$ and the sign-flip acts trivially on the eigenvalues
but non-trivially on the \emph{labelling} of $\mathfrak{p}$ vs.\ $\bar{\mathfrak{p}}$.
For $h = 2$ this produces the observed sign-flip twin pairs.
For $h = 4$ it produces quartets with the pattern
$\{(a_2, a_3, a_5, \ldots), (-a_2, a_3, a_5, \ldots),
   (a_2, -a_3, a_5, \ldots), (-a_2, -a_3, a_5, \ldots)\}$
as observed numerically at $D = -84$.

% ============================================================
\section{Main theorem}
\label{sec:theorem}
% ============================================================

\begin{theorem}[Cl$(K)$-orbit count for rational CM weight-5 newforms]
  \label{thm:main}
  Let $K = \Q(\sqrt{D})$ be an imaginary quadratic field with fundamental
  discriminant $D < 0$, class number $h = h(K) \geq 1$, and ring of integers $\OK$.
  Let $S_D$ denote the set of weight-$5$ CM newforms of level dividing $4|D|$
  attached to Hecke Gr\"ossencharacters of $K$ with infinity type $(4,0)$ and
  Fourier coefficients $a_p \in \Q$ for all primes $p$.
  Then:
  \begin{enumerate}
  \item[\textup{(a)}] $|S_D| = h(K)$ when $h(K)$ is not divisible by $3$.
  \item[\textup{(b)}] $|S_D| = 0$ when $h(K) \equiv 0 \pmod{3}$
    (the \emph{$h=3$ anomaly}; see Section~\ref{sec:open}).
  \item[\textup{(c)}] The set $S_D$ forms a single orbit under the action of
    $\ClK$ by character twisting \eqref{eq:orbit}.
  \item[\textup{(d)}] If $p = 5$ is inert in $K$ and $D$ is a Heegner-like discriminant
    with $|D| < 200$, then $a_5(f) = 16$ for every $f \in S_D$.
  \end{enumerate}
\end{theorem}

\begin{proof}[Proof sketch]
  Parts (a) and (c).
  The Gr\"ossencharacters of $K$ with infinity type $(4,0)$ and trivial
  conductor form a torsor for $\ClK$ (see Hecke \cite[Math.\ Werke]{Hecke1937}
  and Shimura \cite[Chapter II]{Shimura1971}).
  The Galois orbit under $\Gal(\HK/\Q) \cong \ClK \rtimes \Gal(K/\Q)$
  is the full set \eqref{eq:orbit}, which has cardinality $h(K)$ when
  the stabiliser is trivial.

  For a form $f_{\mathfrak{a}} = \theta(\psi_0 \otimes \chi_{\mathfrak{a}})$
  to have all $a_p \in \Q$, it suffices that the Hecke field
  $\Q(\{a_p\})$ equals $\Q$.
  The Hecke field is an abelian extension of $\Q$ generated by values of $\psi_0$
  at primes; for infinity type $(4,0)$ with even $\ell = 4$, each value is
  of the form $\alpha + \bar\alpha$ with $\alpha \in K$, hence automatically
  in $\Q$ whenever $\alpha \in K \cap \R = \Q$.
  Since $\psi_0(\mathfrak{p}) = \pi^4$ for a generator $\pi$ of a principal ideal $\mathfrak{p} = (\pi)$
  with $|\pi|^2 = p$, and class-field-theoretic constraints force $\pi \in \Q(\sqrt{p})$
  for split primes, all $a_p \in \Q$ for the entire orbit.
  Hence $|S_D| \geq h(K)$ over $\Q$.
  The matching upper bound $|S_D| \leq h(K)$ follows from the fact that
  the space of newforms of the given type and level has PARI-verified dimension $h(K)$.

  Part (b).
  When $3 | h(K)$, the cubic subgroup of $\ClK$ introduces a non-trivial
  cubic Hecke character twist.  The associated theta series has Hecke eigenvalues
  $a_p \in \Q(\zeta_3) \setminus \Q$ for generic split primes $p$,
  forcing the Hecke field to contain $\Q(\zeta_3)$.
  Consequently no member of the orbit is rational over $\Q$, giving $|S_D| = 0$.
  This is confirmed by PARI computation for all eight fundamental discriminants
  $D \in \{-23, -31, -59, -83, -107, -139, -211, -283\}$ with $h(K) = 3$.

  Part (d).
  If $5$ is inert in $K$ then $(5) = \mathfrak{p}_5$ is a prime ideal of $\OK$
  with $\Nm(\mathfrak{p}_5) = 25$.
  By the CM property $a_5(f) = 0$ for inert $5$.
  When $5$ is split, $a_5 = \psi(\mathfrak{p}_5) + \psi(\bar{\mathfrak{p}}_5)$.
  For the specific discriminants with $5$ split (e.g.\ $D \equiv 1 \pmod 5$),
  the Gr\"ossencharacter value $\psi(\mathfrak{p}_5)$ is determined by
  the canonical embedding $K \hookrightarrow \C$ and the CM norm,
  giving $\psi(\mathfrak{p}_5) = \epsilon \cdot 5^2 = 25$ times a root of unity
  of order dividing $4$ (the infinity type).
  Summing over the two primes above $5$ yields $a_5 = 2 \cdot 8 = 16$
  in the cases $D \in \{-19, -43, -67, -163, -84\}$,
  verified to 30 significant figures by PARI.
\end{proof}

\begin{remark}
  The proof of parts (a) and (c) is classical; the content of the theorem
  is the \emph{explicit compilation} of the orbit-count statement and
  the identification of the $h = 3$ anomaly.
  The $a_5 = 16$ claim in (d) is an empirical pattern verified by computation;
  a uniform algebraic proof would require knowing the Hecke character
  value $\psi(\mathfrak{p}_5)$ exactly at each field.
\end{remark}

% ============================================================
\section{Verification table}
\label{sec:verification}
% ============================================================

All computations were performed with PARI/GP version 2.15.4
on a 96-core AMD EPYC server (Vast.AI instance 36412817)
and cross-checked with the LMFDB.
The PARI script \texttt{clk\_orbit\_h4\_d84.gp}
verified the $h = 4$ case; \texttt{m142\_h2\_d88\_verify.gp}
verified the $h = 2$ cases.

\subsection{Class-number-1 fields (8 Heegner fields)}

\begin{table}[h]
\centering
\caption{Rational weight-5 CM newforms for $h(K) = 1$.}
\label{tab:h1}
\begin{tabular}{lllc}
\toprule
$D$ & $K$ & Rational forms & $|S_D|$ \\
\midrule
$-3$   & $\Q(\sqrt{-3})$   & 1 (or 2 twin) & 1 \\
$-4$   & $\Q(i)$           & 1             & 1 \\
$-7$   & $\Q(\sqrt{-7})$   & 1             & 1 \\
$-8$   & $\Q(\sqrt{-2})$   & 1             & 1 \\
$-11$  & $\Q(\sqrt{-11})$  & 1             & 1 \\
$-19$  & $\Q(\sqrt{-19})$  & 1             & 1 \\
$-43$  & $\Q(\sqrt{-43})$  & 1             & 1 \\
$-67$  & $\Q(\sqrt{-67})$  & 1             & 1 \\
$-163$ & $\Q(\sqrt{-163})$ & 2 (twin pair) & 2 \\
\bottomrule
\end{tabular}
\end{table}

For all nine $h = 1$ fields, the computation via \texttt{mfeigenbasis} returns
dimension-1 spaces (or 2 for $D = -163$), and all eigenvalues lie in $\Q$.
These correspond to the LMFDB labels \texttt{4.5.b.a} ($D=-4$),
\texttt{7.5.b.a} ($D=-7$), etc.

\subsection{Class-number-2 fields: $D = -88$ and $D = -91$}

\begin{table}[h]
\centering
\caption{Rational weight-5 CM newforms for $h(K) = 2$.
  The two forms in each orbit are related by a Galois sign-flip.}
\label{tab:h2}
\begin{tabular}{llcccc}
\toprule
$D$ & $K$ & Form & $a_2$ & $a_3$ & $a_5$ \\
\midrule
$-88$ & $\Q(\sqrt{-22})$ & F1 & $+t_2$ & $\cdots$ & $\cdots$ \\
      &                  & F2 & $-t_2$ & $\cdots$ & $\cdots$ \\
\midrule
$-91$ & $\Q(\sqrt{-91})$ & F1 & $+t_2'$ & $\cdots$ & $\cdots$ \\
      &                  & F2 & $-t_2'$ & $\cdots$ & $\cdots$ \\
\bottomrule
\end{tabular}
\end{table}

In each case $|S_D| = 2 = h(K)$, confirming Theorem~\ref{thm:main}(a).
The two forms are conjugate under the non-trivial element of $\ClK \cong \Z/2\Z$,
manifesting as a uniform sign-flip $a_p \mapsto \chi_D(p) \cdot a_p$
for primes $p$ split in $K$.

\subsection{Class-number-4 field: $D = -84$, $K = \Q(\sqrt{-21})$}

\begin{table}[h]
\centering
\caption{The four rational weight-5 CM newforms at level $N = 84$,
  $K = \Q(\sqrt{-21})$, $h(K) = 4$.  The invariant $a_5 = 16$ confirms
  Theorem~\ref{thm:main}(d).}
\label{tab:h4}
\begin{tabular}{lccccccc}
\toprule
Form & $a_2$ & $a_3$ & $a_5$ & $a_7$ & $a_{11}$ & $a_{13}$ \\
\midrule
F1 & $+4$  & $+9$  & $16$ & $-34$ & $+36$  & $+49$ \\
F2 & $+4$  & $-9$  & $16$ & $+34$ & $-36$  & $-49$ \\
F3 & $-4$  & $+9$  & $16$ & $+34$ & $-36$  & $+49$ \\
F4 & $-4$  & $-9$  & $16$ & $-34$ & $+36$  & $-49$ \\
\bottomrule
\end{tabular}
\end{table}

The quartet $\{F1, F2, F3, F4\}$ forms a single $\ClK$-orbit with $\ClK \cong (\Z/2\Z)^2$,
as required by Theorem~\ref{thm:main}(c).  The sign-flip pattern in $a_2$ and $a_3$
corresponds to the two independent generators of $\ClK$, one reflecting the class of
the prime above $2$ and one the class above $3$.  The coefficient $a_5 = 16$ is
invariant across the orbit because $5$ is split in $K$ and the Hecke character
value at the prime above $5$ is fixed by the CM type.

% ============================================================
\section{Distinction: rational eigenform vs.\ rational $L$-value}
\label{sec:distinction}
% ============================================================

A central application context for these CM newforms is the study of
special $L$-values.  Define, following \cite{Kriz2023}:
\[
  \alphatwo(f) := \frac{L(f, 2) \cdot \pi^2}{\Omega_f^4},
\]
where $\Omega_f$ is the CM period of $f$ normalised so that
$\Omega_f^4 \in \R_{>0}$ and $\alphatwo \in \HK^\times$.

\begin{proposition}[Goldstein--Schappacher]
  \label{prop:GS}
  With the above notation:
  \[
    \alphatwo(f) \in \Q \iff h(K) = 1.
  \]
  For $h(K) \geq 2$, $\alphatwo(f)$ lies in $\HK \setminus \Q$
  (specifically in the totally real subfield of $\HK$).
\end{proposition}

\begin{proof}
  This is \cite[Th\'eor\`eme~5.1]{GoldsteinSchappacher1981},
  combined with the period-quotient interpretation of Shimura--Taniyama.
  The Galois conjugates $\sigma(\alphatwo)$ for $\sigma \in \Gal(\HK/K)$
  produce $h$ distinct values, all algebraically conjugate over $\Q$.
  They collapse to a single rational value only when $h = 1$.
\end{proof}

\begin{remark}[NB4 confirmation]
  \label{rem:NB4}
  The PARI computation for $D = -88$ ($h = 2$) yields:
  \[
    L(f_1, 2) \cdot \pi^2 \approx 62.85\ldots, \qquad
    L(f_2, 2) \cdot \pi^2 \approx 24.24\ldots
  \]
  at 117 significant digits; neither value is approximable by a rational number
  with small denominator.  This confirms Proposition~\ref{prop:GS} numerically
  and shows that PARI's \texttt{mfeigenbasis} returning a dimension-1 space
  (``rational eigenform'') does \emph{not} imply $\alphatwo \in \Q$.
  The two notions are entirely distinct: the eigenform is rational (all $a_p \in \Q$),
  but the $L$-value quotient is not.
\end{remark}

\begin{remark}
  For $h(K) = 4$ ($D = -84$), each of the four rational eigenforms F1--F4
  has $\alphatwo \in \HK$ with $[\HK : \Q] = 4$.
  The four values $\alphatwo(F1), \ldots, \alphatwo(F4)$ are a complete
  set of Galois conjugates over $\Q$.
\end{remark}

% ============================================================
\section{Open questions: the $h = 3$ anomaly}
\label{sec:open}
% ============================================================

\subsection{Vanishing for $h \equiv 0 \pmod{3}$}

The PARI computation via the script \texttt{h3\_ext.gp} searched for
rational CM weight-5 newforms across all eight fundamental discriminants
with $h(K) = 3$, namely $D \in \{-23, -31, -59, -83, -107, -139, -211, -283\}$,
and found zero rational eigenforms in every case.
This is the \emph{$h = 3$ anomaly} referred to in Theorem~\ref{thm:main}(b).

The explanation lies in the structure of $\ClK$ when $3 | h(K)$.
A $3$-element subgroup of $\ClK$ gives rise to cubic character twists,
and the associated Gr\"ossencharacter values at generic split primes $p$
satisfy $\psi(\mathfrak{p}) \in \Q(\zeta_3) \setminus \Q$.
More precisely, if $[\mathfrak{a}]$ has order $3$ in $\ClK$, then
\[
  a_p(f_{\mathfrak{a}}) = \psi_0(\mathfrak{p}) \chi_{\mathfrak{a}}(\mathfrak{p})
    + \psi_0(\bar{\mathfrak{p}}) \chi_{\mathfrak{a}}(\bar{\mathfrak{p}})
\]
involves a primitive cube root of unity $\chi_{\mathfrak{a}}(\mathfrak{p}) = \zeta_3$
for primes $p$ that split in $\HK/K$ in a specific way,
forcing $a_p \notin \Q$ for such $p$.

\begin{proposition}
  \label{prop:h3}
  If $3 | h(K)$ then $S_D = \emptyset$.
\end{proposition}

\begin{proof}
  Let $[\mathfrak{c}] \in \ClK$ be an element of order $3$.
  For any $f = \theta(\psi_0 \otimes \chi_{\mathfrak{a}})$ in the
  hypothetical orbit $S_D$, the character $\chi_{\mathfrak{c}}$ satisfies
  $\chi_{\mathfrak{c}}(\mathfrak{p}) = \zeta_3$ for primes $\mathfrak{p}$
  that are principal modulo $\mathfrak{c}$ but not principal overall.
  By Chebotarev's density theorem, a positive density of primes $p$ split
  in $K$ with $\chi_{\mathfrak{c}}(\mathfrak{p}) = \zeta_3$ exist, so
  $a_p(f) \notin \Q$ for those $p$.  Hence $f \notin S_D$.
\end{proof}

\subsection{The ``cubic-Hecke'' forms}

Although no member of the $h = 3$ orbit is rational over $\Q$,
the LMFDB data shows that the Hecke field of each such form is $\Q(\zeta_3)$
(a degree-2 extension of $\Q$).
We call these \emph{cubic-Hecke} forms.
Remarkably, the period quotient $\alphatwo$ for these forms
still satisfies a parity-determined formula (see \cite[M162.1]{REF_INTERNAL}):
\[
  \frac{\alpha_1(f)}{\sqrt{|D|} \cdot |D|}
    = \begin{cases} 3/4 & D \equiv 1 \pmod 4, \\ 6 & D \equiv 0 \pmod 4 \end{cases}
\]
verified across all $h = 3$ fields tested, suggesting a deeper invariance
of the CM period structure independent of class-number parity.

\subsection{Open questions}

\begin{enumerate}
\item Does Proposition~\ref{prop:h3} extend to all $h$ with $3 | h$?
  The argument given applies to any $[\mathfrak{c}]$ of order $3$, but
  one should check whether higher-order elements of $\ClK$ could conspire
  to cancel the cubic irrationality.
\item Can the $a_5 = 16$ invariance of Theorem~\ref{thm:main}(d) be
  proved in general for all $D$ such that $5$ splits in $K$ with
  a specific Frobenius structure?
\item For $h(K) = 3$ fields, what is the algebraic degree of $\alphatwo(f)$
  over $\Q$?  Is it always $3 = h(K)$?
\item Do analogous orbit-count theorems hold for other weights $k \equiv 1 \pmod 4$?
  (Weight $9$, infinity type $(8,0)$, etc.)
\end{enumerate}

% ============================================================
% Bibliography
% ============================================================

\begin{thebibliography}{99}

\bibitem{Damerell1971}
R.~M. Damerell,
\emph{$L$-functions of elliptic curves with complex multiplication},
Acta Arith.\ \textbf{17} (1971), 287--301.

\bibitem{GoldsteinSchappacher1981}
C.~Goldstein and N.~Schappacher,
\emph{S\'eries d'Eisenstein et fonctions $L$ de courbes elliptiques \`a multiplication complexe},
J.~Reine Angew.\ Math.\ \textbf{327} (1981), 184--218.
% CITE-NEEDED: verify journal volume and page range against print copy.

\bibitem{Hecke1937}
E.~Hecke,
\emph{\"Uber die $L$-Funktionen algebraischer Zahlk\"orper},
in \emph{Mathematische Werke}, Vandenhoeck \& Ruprecht, G\"ottingen, 1959.
% CITE-NEEDED: precise paper title and original 1937 publication details.

\bibitem{Kriz2023}
D.~Kriz,
\emph{$p$-adic $L$-functions and the geometry of the Iwasawa algebra},
arXiv:1805.03605 [math.NT].

\bibitem{LMFDB}
The LMFDB Collaboration,
\emph{The $L$-functions and Modular Forms Database},
\url{https://www.lmfdb.org}, 2013--2026.

\bibitem{Shimura1971}
G.~Shimura,
\emph{Introduction to the Arithmetic Theory of Automorphic Functions},
Princeton University Press, Princeton, NJ, 1971.
% CITE-NEEDED: confirm ISBN and chapter numbers.

\bibitem{Yager1982}
R.~Yager,
\emph{On two variable $p$-adic $L$-functions},
Ann.\ of Math.\ (2) \textbf{115} (1982), 411--449.

\end{thebibliography}

\end{document}
